\section*{第 \ref{ch:g9:natural-risks} 章——与自然风险共存}

\begin{solution}{exo:g9:natural-risks:1}
(a) 危害。(b) \omterm{def:g9:natural-risks:vulnerability}{脆弱性}。(c) \omterm{def:g9:natural-risks:vulnerability}{脆弱性}。(d) 危害。(e) 两者兼有：
泥流是危害；把城镇建在泥流路径上，使其
变得脆弱。
\end{solution}

\begin{solution}{exo:g9:natural-risks:2}
(a) 错误：无人居住之处就无人受害，因此\omterm{def:g9:natural-risks:risk}{风险}较低。
(b) 正确。(c) 正确。(d) 错误：预警并不能改变\omterm{def:g4:volcanoes-earthquakes-first-look:volcano}{火山}；它
降低了人群的\omterm{def:g9:natural-risks:vulnerability}{脆弱性}，使人们能够及时撤离。
\end{solution}

\begin{solution}{exo:g9:natural-risks:3}
(a) 能，提前数天。(b) 能，提前数小时至数天。(c) 否：\omterm{def:g4:volcanoes-earthquakes-first-look:earthquake}{地震}
无法预测。(d) 能：\omterm{def:g7:earthquakes:tsunami}{海啸}一旦发生，穿越\omterm{def:g2:water-on-earth:ocean}{海洋}需要
数小时，而预警传播得更快。
\end{solution}

\begin{solution}{exo:g9:natural-risks:4}
(a) 地面的震动：\omterm{def:g7:volcanoes:magma}{岩浆}上升时出现成群的小
\omterm{def:g4:volcanoes-earthquakes-first-look:earthquake}{地震}。(b) 斜坡倾斜：\omterm{def:g4:volcanoes-earthquakes-first-look:volcano}{火山}隆起。(c) 地面如何
移动，以毫米计：\omterm{def:g7:volcanoes:magma}{岩浆}聚集时的隆起。(d) 烟羽
中的气体：\omterm{def:g7:volcanoes:magma}{岩浆}上升时逸出的气体增多。
\end{solution}

\begin{solution}{exo:g9:natural-risks:5}
(a) $45 \div 2 = \qty{22.5}{km/h}$。(b) $4 \times 2 = \qty{8}{km}$，足以
到达高地。(c) 无人目睹\omterm{def:g4:volcanoes-earthquakes-first-look:eruption}{喷发}，\omterm{def:g4:volcanoes-earthquakes-first-look:eruption}{喷发}被\omterm{def:g5:natural-hazards:storm}{风暴}遮挡，
也无人向他们发出预警：当时没有快速向城镇
传递信息的途径，也没有疏散计划。(d) 例如：
一份带有通往高地路线并预先演练的疏散计划；
一条从科学家到城镇的预警链；在\omterm{def:g4:volcanoes-earthquakes-first-look:volcano}{火山}上增设
监测仪器。
\end{solution}

\begin{solution}{exo:g9:natural-risks:6}
(a) \omterm{def:g5:water-shapes-the-land:flood}{洪水}高出\omterm{def:g1:landscapes-around-us:water}{河流} \qty{10}{m}。(b) 约 \qty{0.47}{km}（\qty{2}{m} \omterm{def:g5:water-shapes-the-land:flood}{洪水}），
\qty{0.56}{km}（\qty{5}{m} \omterm{def:g5:water-shapes-the-land:flood}{洪水}）以及
\qty{0.67}{km}（\qty{10}{m} \omterm{def:g5:water-shapes-the-land:flood}{洪水}）（在约
\qty{0.05}{km} 之内的读数均可）。(c) 只有 \qty{10}{m} \omterm{def:g5:water-shapes-the-land:flood}{洪水}：该处地面
高于\omterm{def:g1:landscapes-around-us:water}{河流} 2 和 \qty{5}{m}。(d) \omterm{def:g1:landscapes-around-us:water}{河流}只需
小幅上涨即可到达该处，而小\omterm{def:g5:water-shapes-the-land:flood}{洪水}比大\omterm{def:g5:water-shapes-the-land:flood}{洪水}
更频繁地出现。
\end{solution}

\begin{solution}{exo:g9:natural-risks:7}
(a) 危害高、\omterm{def:g9:natural-risks:vulnerability}{脆弱性}低：中等\omterm{def:g9:natural-risks:risk}{风险}。(b) 危害低、
\omterm{def:g9:natural-risks:vulnerability}{脆弱性}高：中等\omterm{def:g9:natural-risks:risk}{风险}。(c) 危害高、\omterm{def:g9:natural-risks:vulnerability}{脆弱性}高：高
\omterm{def:g9:natural-risks:risk}{风险}。(d) 危害低、\omterm{def:g9:natural-risks:vulnerability}{脆弱性}低：低\omterm{def:g9:natural-risks:risk}{风险}。
\end{solution}

\begin{solution}{exo:g9:natural-risks:8}
(a) 依次为 C、B、A。(b) 支撑杆构成三角形，三角形
在不拉伸或压缩支撑杆的情况下无法倾斜，框架因此保持形状，
不会歪倒。(c) 隔震器随地面
在其下方前后移动而让位，因此传到学校的运动
大为减少。
\end{solution}

\begin{solution}{exo:g9:natural-risks:9}
(a) 每年 $1 \div 100 = 1~\%$。(b) 每年的概率都相同，
与前一年发生什么无关：\omterm{def:g5:water-shapes-the-land:flood}{洪水}明年仍可能再次发生。
(c) $300 \div 100 = 3$，平均为 3 次\omterm{def:g5:water-shapes-the-land:flood}{洪水}。
\end{solution}

\begin{solution}{exo:g9:natural-risks:10}
(a) $3600 \div 15 = 240$ 个读数。(b) $240 \times 24 \times 365 =
\num{2102400} \approx 2.1 \times 10^6$ 个读数。(c) $2400 \div 800 =
\qty{3}{h}$。(d) 在这三小时内，预警中心可以向
\omterm{def:g1:landscapes-around-us:water}{海岸}发出警告，人们也可以转移到高地。
\end{solution}

\begin{solution}{exo:g9:natural-risks:11}
(a) $4 \times 3 \times 3.8 = \qty{45.6}{L}$。(b) $45.6 \div 1.5 = 30.4$，
因此需要 31 瓶。(c) 例如：耐储存的食物，以便在商店关门时食用；
手电筒，以备停电时使用；靠电池或手摇供电的收音机，
用来接收预警；急救包，用来处理
轻微伤害；保暖毯子。(d) 以便在家人必须
迅速撤离时，能够立即带走。
\end{solution}

\begin{solution}{exo:g9:natural-risks:12}
例如：(1) 逐个\omterm{def:g1:landscapes-around-us:landform}{山谷}绘制泥流的\omterm{def:g9:natural-risks:hazard-map}{危害分布图}
（危害方面）；(2) 在\omterm{def:g4:volcanoes-earthquakes-first-look:volcano}{火山}上安装\omterm{def:g7:earthquakes:seismometer}{地震仪}、倾斜仪、GPS 接收器和气体
监测仪器（危害方面）；(3) 建立快速预警链，
从观测站到城镇，配以警报器（\omterm{def:g9:natural-risks:vulnerability}{脆弱性}
方面）；(4) 规划并标明通往高地的路线，并
以演习加以练习（\omterm{def:g9:natural-risks:vulnerability}{脆弱性}方面）；(5) 不在
泥流路径上新建住房或学校（\omterm{def:g9:natural-risks:vulnerability}{脆弱性}方面）。
\end{solution}

\begin{solution}{pb:g9:natural-risks:1}
\textbf{1.} \omterm{def:g1:landscapes-around-us:water}{河流}在 C 旁的高程为 \qty{143}{m}，在 A 旁为 \qty{137}{m}，
在 B 旁为 \qty{126}{m}；图上 C 位于西侧，B 位于东侧，
因此\omterm{def:g1:landscapes-around-us:water}{河流}向东流。

\textbf{2.} A：$137 - 137 = \qty{0}{m}$；B：$276 - 126 = \qty{150}{m}$；
C：$236 - 143 = \qty{93}{m}$。

\textbf{3.} 只有 A，它比\omterm{def:g1:landscapes-around-us:water}{河流}高出不到 \qty{5}{m}。

\textbf{4.} 约 \qty{0.57}{km}（介于 0.5 与 \qty{0.65}{km} 之间）。

\textbf{5.} 在 \qty{100}{m} 范围内的上升量，是
\qty{200}{m} 范围内上升量的一半：A $1.4 \div 2 = 0.7~\%$；B $83 \div 2 = 41.5~\%$；
C $20 \div 2 = 10~\%$。

\textbf{6.} B：其坡度很陡，松散的\omterm{def:g3:stones-pebbles-sand:rock}{岩石}和土壤
容易沿陡坡滑落，尤其是在大雨之后。

\textbf{7.} A：约 \qty{1}{km}；B：几乎位于\omterm{def:g7:earthquakes:fault}{断层}上（小于
\qty{0.05}{km}）；C：约 \qty{1.3}{km}。

\textbf{8.} 若\omterm{def:g7:earthquakes:fault}{断层}发生移动，地面会沿\omterm{def:g7:earthquakes:fault}{断层}断裂并错动，
这可能把建筑撕裂；并且，位于\omterm{def:g7:earthquakes:fault}{断层}上的\omterm{def:g4:volcanoes-earthquakes-first-look:earthquake}{地震}会使该处地面
剧烈震动。

\textbf{9.}

\begin{center}
{\small
\begin{tabular}{|c|c|c|c|}
\hline
场地 & 位于洪泛区？ & 陡坡？ & 位于\omterm{def:g7:earthquakes:fault}{断层}上？ \\
\hline
A & 是 & 否 & 否 \\
B & 否 & 是 & 是 \\
C & 否 & 否 & 否 \\
\hline
\end{tabular}}
\end{center}

\textbf{10.} C 场地：它是唯一三项均为“否”的场地，高出河面 \qty{93}{m}，
坡度平缓，为 \qty{10}{\%}，
距\omterm{def:g7:earthquakes:fault}{断层}超过 \qty{1}{km}。

\textbf{11.} 例如：框架中设置斜撑和剪力墙，
采用牢固的节点，或在建筑下方安装隔震器；学生进行
\omterm{def:g4:volcanoes-earthquakes-first-look:earthquake}{地震}演练（伏地、掩护、抓牢）。

\textbf{12.} 28 座房屋中，\textbf{17 座}位于洪泛区内：
$17 \div 28 \approx 0.61$，约为 5 座中有 3 座。
\end{solution}
